\newtoggle{withbonus} \toggletrue{withbonus} % set true to show bonus points in grade table

% setting underlying course name (Grundlagenfach Mathematik, §-Unterricht, \ldots)
\newcommand{\mycourse}{Grundlagenfach Informatik}
% set date
\newcommand{\mydate}{Dienstag, 07.04.2026}

% name of the class
\newcommand{\myclass}{Klasse 2c}
\title{\textbf{Hardware, Endliche Automaten, Laufzeitkomplexität}}
\author{Thomas Graf\\
	\mycourse, \myclass}
\date{\mydate}
\maketitle
\thispagestyle{headandfoot}

\begin{center}
	\fbox{\parbox{140mm}{\centering
			\begin{itemize}
				\item Dauer der Prüfung: $45$ Minuten
				      %\item \textbf{Zeige in jeder Aufgabe unbedingt Deine Schritte!}
				\item Erlaubte Hilfsmittel:
				      \begin{itemize}
					      \item Schreibzeug
					            %\item Taschenrechner
				      \end{itemize}
			\end{itemize}}}
\end{center}
\vspace{10mm}

\begin{center}
	Vorname: \rule{45mm}{0.8pt}\qquad Nachname: \rule{45mm}{0.8pt}
\end{center}
\vspace{20mm}

\begin{center}
	\iftoggle{withbonus} {
		\combinedgradetable[h][questions]
	} {
		\gradetable[h][questions]
	}
\end{center}
\vspace{15mm}
\begin{center}
	Note:
	\vspace{10mm}

	\rule{8mm}{1.2pt}
\end{center}
%\thispagestyle{empty}
\clearpage


\begin{questions}
	\question
	\begin{parts}
		\part[1] Welche logische Funktion wird durch zwei parallel geschaltete Schalter $A$ und $B$ realisiert, damit eine LED leuchtet?
		\begin{checkboxes}
			\choice \textsc{XOR}
			\correctchoice \textsc{OR}
			\choice \textsc{AND}
		\end{checkboxes}

		\part[1] Welche Formel beschreibt das Carry-Out-Bit $C_{out}$ eines Volladdierers?
		\begin{checkboxes}
			\choice $C_{out} = A \oplus B \oplus C_{in}$
			\correctchoice $C_{out} = (A\;\operatorname{AND}\;B)\;\operatorname{OR}\;(C_{in}\;\operatorname{AND}\;(A\oplus B))$
			\choice $C_{out} = A \oplus B$
			\choice keine der obigen Antworten ist richtig
		\end{checkboxes}
	\end{parts}
	\droptotalpoints


	\question[2]{\textbf{Anzahl der Elemente in einer Sprache bestimmen}}

	Es sei $\Sigma = \lrc{0,1,2}$ ein Alphabet und $n\in\N$ gegeben. Bestimmen Sie die Anzahl der verschiedenen Wörter der Länge $n$ über $\Sigma$, welche mindestens eine $0$ und mindestens eine $2$ enthalten.
	\begin{solution}
		Insgesamt gibt es $3^n$ Wörter der Länge $n$ über $\Sigma$.

		Wir zählen die Wörter, welche nicht beide Bedingungen erfüllen. Wörter ohne $0$ gibt es $2^n$, Wörter ohne $2$ gibt es ebenfalls $2^n$. Das Wort, welches weder $0$ noch $2$ enthält, ist nur das Wort $11\ldots 1$ und wurde doppelt gezählt.

		Daher ist die gesuchte Anzahl
		\[
			3^n - 2\cdot 2^n + 1
			=
			3^n - 2^{n+1} + 1.
		\]
	\end{solution}
	\droptotalpoints

	\vspace{5cm}

	\clearpage

	\question{\textbf{Endlichen Automaten zeichnen}}

	Gegeben ist die Sprache
	\[
		L_1 := \setcm{x00y1}{x,y\in\lrc{0,1}^*}.
	\]
	\begin{parts}
		\part[1] Nennen Sie ein kürzestes Wort in $L_1$.
		\begin{solution}
			Ein kürzestes Wort ist
			\[
				001.
			\]
		\end{solution}

		\vspace{4cm}

		\part[2] Entwerfen Sie einen endlichen Automaten in Diagrammform, welcher $L_1$ akzeptiert.
		\begin{solution}
			\begin{figure}[H]
				\centering
				\begin{tikzpicture}[->,>=stealth',shorten >=1pt,node distance=3.2cm,semithick]
					\tikzstyle{every state}=[fill=green!10,draw=black,text=black,minimum size=1.4cm]

					\node[initial,state] (q0) {$q_0$};
					\node[state] (q1) [right of=q0] {$q_1$};
					\node[state] (q2) [right of=q1] {$q_2$};
					\node[state] (q3) [right of=q2] {$q_3$};
					\node[state,accepting] (q4) [right of=q3] {$q_4$};

\path
(q0) edge [loop above] node {$1$} (q0)
(q0) edge [above] node {$0$} (q1)
(q1) edge [above] node {$0$} (q2)
(q1) edge [bend left, above] node {$1$} (q0)
(q2) edge [loop above] node {$0$} (q2)
(q2) edge [above] node {$1$} (q4)
(q4) edge [loop above] node {$1$} (q4)
(q4) edge [bend left, below] node {$0$} (q3)
(q3) edge [loop above] node {$0$} (q3)
(q3) edge [bend left, above] node {$1$} (q4);
				\end{tikzpicture}
				\caption{Endlicher Automat für die Sprache $L_1$}
			\end{figure}
		\end{solution}
	\end{parts}
	\droptotalpoints

\clearpage
\question{\textbf{\red{Nicht reguläre} Sprache}}

Gegeben sei die Sprache
\begin{align*}
	L_2 := \setcm{w\in\lrc{a,b}^*}{\abs{w}_a = 2\abs{w}_b}
\end{align*}
über dem Alphabet $\lrc{a,b}$.

Dabei bezeichnet $\abs{w}_a$ die Anzahl der Buchstaben $a$ im Wort $w$
und $\abs{w}_b$ die Anzahl der Buchstaben $b$ im Wort $w$.

\begin{parts}

\part[1]
Nenne genau drei verschiedene Wörter, die in $L_2$ liegen.

\begin{solution}
\[
aab,\qquad aba,\qquad baa
\]
\end{solution}

\vspace{3cm}

\part[2]
Zeige, dass $L_2$ \red{nicht regulär} ist.

\begin{solution}
Angenommen, $L_2$ sei regulär. Dann existiert ein endlicher Automat
\[
A=(Q,\lrc{a,b},\delta,q_0,F),
\]
welcher $L_2$ akzeptiert.

Sei $m:=|Q|$ die Anzahl seiner Zustände.

Wir betrachten nun die $m+1$ verschiedenen Wörter
\[
a^2,a^4,a^6,\ldots,a^{2m+2}.
\]

Nach dem Schubfachprinzip existieren zwei Zahlen
\[
1\le i<j\le m+1
\]
mit
\[
\hat{\delta}(q_0,a^{2i})
=
\hat{\delta}(q_0,a^{2j}).
\]

Nach der im Unterricht bewiesenen Aussage gilt deshalb für jedes Wort
$z\in\lrc{a,b}^*$
\[
a^{2i}z\in L_2
\iff
a^{2j}z\in L_2.
\]

Wir wählen nun
\[
z:=b^i.
\]

Dann gilt
\[
a^{2i}b^i\in L_2,
\]
da
\[
\abs{a^{2i}b^i}_a=2i
\qquad\text{und}\qquad
\abs{a^{2i}b^i}_b=i.
\]

Andererseits gilt wegen $j>i$
\[
a^{2j}b^i\notin L_2,
\]
da
\[
\abs{a^{2j}b^i}_a=2j>2i=2\abs{a^{2j}b^i}_b.
\]

Dies ist ein Widerspruch.

Also war unsere Annahme falsch und $L_2$ ist nicht regulär.
\end{solution}

\end{parts}

\droptotalpoints
	\clearpage

	\question Bestimme jeweils die Worst-Case-Laufzeit der folgenden Algorithmen in Abhängigkeit von $n\in\N$.
	\begin{parts}
		\part[1] \leavevmode
		\begin{lstlisting}[language=Python]
result = 0
for i in range(n):
    for j in range(i):
        result += 1
		\end{lstlisting}
		\begin{checkboxes}
			\choice linear
			\choice proportional zu $\sqrt{n}$
			\choice logarithmisch
			\correctchoice quadratisch
			\choice proportional zu $n\log(n)$
			\choice keine der obigen Antworten ist richtig
		\end{checkboxes}

		\part[1] \leavevmode
		\begin{lstlisting}[language=Python]
result = 0
i = 1
while i * i <= n:
    result += i
    i += 1
		\end{lstlisting}
		\begin{checkboxes}
			\choice linear
			\correctchoice proportional zu $\sqrt{n}$
			\choice logarithmisch
			\choice quadratisch
			\choice proportional zu $n\log(n)$
			\choice keine der obigen Antworten ist richtig
		\end{checkboxes}

		\part[1] \leavevmode
		\begin{lstlisting}[language=Python]
result = 0
for i in range(n):
    j = 1
    while j <= n:
        result += 1
        j *= 2
		\end{lstlisting}
		\begin{checkboxes}
			\choice linear
			\choice proportional zu $\sqrt{n}$
			\choice logarithmisch
			\choice quadratisch
			\correctchoice proportional zu $n\log(n)$
			\choice keine der obigen Antworten ist richtig
		\end{checkboxes}

		\part[1] \leavevmode
		\begin{lstlisting}[language=Python]
result = 0
i = n
while i >= 1:
    result += 1
    i //= 3
		\end{lstlisting}
		\begin{checkboxes}
			\choice proportional zu $n\log(n)$
			\choice proportional zu $\sqrt{n}$
			\correctchoice logarithmisch
			\choice quadratisch
			\choice linear
			\choice keine der obigen Antworten ist richtig
		\end{checkboxes}
	\end{parts}
	\droptotalpoints

\end{questions}